
\documentclass{article}
\usepackage{colortbl}
\usepackage{makecell}
\usepackage{multirow}
\usepackage{supertabular}

\begin{document}

\newcounter{utterance}

\centering \large Interaction Transcript for game `cladder', experiment `full\_v1.5\_default', episode 2591 with Qwen3.5{-}9B.
\vspace{24pt}

{ \footnotesize  \setcounter{utterance}{1}
\setlength{\tabcolsep}{0pt}
\begin{supertabular}{c@{$\;$}|p{.15\linewidth}@{}p{.15\linewidth}p{.15\linewidth}p{.15\linewidth}p{.15\linewidth}p{.15\linewidth}}
    \# & \multicolumn{2}{c}{Player} && \multicolumn{2}{c}{Game Master} \\
    \hline

    \theutterance \stepcounter{utterance}  
    & & & \multicolumn{4}{p{0.6\linewidth}}{
        \cellcolor[rgb]{0.9,0.9,0.9}{
            \makecell[{{p{\linewidth}}}]{
                \texttt{\tiny{[P1$\langle$GM]}}
                \texttt{You are an expert in causal inference. The following question is not a typical commonsense query, but rather a meticulously designed question created by a professor specializing in causal inference, intended to assess the students' mastery of the course content.} \\
\texttt{Imagine a self{-}contained, hypothetical world with only the following conditions, and without any unmentioned factors or causal relationships: Gender has a direct effect on smoking and lung cancer. Smoking has a direct effect on tar deposit. Tar deposit has a direct effect on lung cancer. Gender is unobserved. For nonsmokers, the probability of high tar deposit is 77\%. For smokers, the probability of high tar deposit is 26\%. For nonsmokers and with no tar deposit, the probability of lung cancer is 71\%. For nonsmokers and with high tar deposit, the probability of lung cancer is 45\%. For smokers and with no tar deposit, the probability of lung cancer is 63\%. For smokers and with high tar deposit, the probability of lung cancer is 37\%. The overall probability of smoking is 39\%. Will smoking decrease the chance of lung cancer?} \\
\texttt{Start your answer with 'Yes' or 'No', followed by additional reasoning or evidence to support your explanation.} \\
            }
        }
    }
     \\ \\

    \theutterance \stepcounter{utterance}  
    & \multicolumn{4}{p{0.6\linewidth}}{
        \cellcolor[rgb]{0.561,0.737,0.561}{
            \makecell[{{p{\linewidth}}}]{
                \texttt{\tiny{[P1$\rangle$GM]}}
                \texttt{No} \\
\\ 
\texttt{**Reasoning:**} \\
\\ 
\texttt{To determine if smoking decreases the chance of lung cancer, we must calculate the causal effect of smoking on lung cancer, denoted as \$P(LungCancer | do(Smoking))\$. In this hypothetical world, the causal structure is defined as:} \\
\texttt{\$Gender \textbackslash{}rightarrow Smoking\$} \\
\texttt{\$Smoking \textbackslash{}rightarrow TarDeposit\$} \\
\texttt{\$TarDeposit \textbackslash{}rightarrow LungCancer\$} \\
\\ 
\texttt{Since \$Gender\$ is unobserved, we cannot simply average over it using standard adjustment formulas without knowing the distribution of smoking given gender. However, we can analyze the mechanism directly. The probability of lung cancer is determined entirely by the state of tar deposits, which are determined by smoking status. Therefore, the causal effect of smoking is the difference in the probability of lung cancer between the world where we force everyone to smoke and the world where we force everyone to not smoke, marginalizing over the distribution of tar deposits induced by those actions.} \\
\\ 
\texttt{Let's calculate the probability of lung cancer under the intervention \$do(Smoking = Yes)\$:} \\
\texttt{1.  **Intervention:** Everyone smokes.} \\
\texttt{2.  **Tar Deposit Distribution:** Since everyone smokes, the probability of high tar deposit is the probability for smokers, which is given as \$P(Tar=High | Smoking) = 0.26\$. Consequently, \$P(Tar=No | Smoking) = 1 {-} 0.26 = 0.74\$.} \\
\texttt{3.  **Lung Cancer Probability:**} \\
\texttt{    *   If Tar is High: \$P(LungCancer | Tar=High) = 0.37\$ (given in the prompt for smokers/high tar).} \\
\texttt{    *   If Tar is No: \$P(LungCancer | Tar=No) = 0.63\$ (given in the prompt for smokers/no tar).} \\
\texttt{    *   \$P(LungCancer | do(Smoking=Yes)) = (0.26 \textbackslash{}times 0.37) + (0.74 \textbackslash{}times 0.63)\$} \\
\texttt{    *   Calculation: \$0.0962 + 0.4662 = 0.5624\$ (or 56.24\%).} \\
\\ 
\texttt{Now, let's calculate the probability of lung cancer under the intervention \$do(Smoking = No)\$:} \\
\texttt{1.  **Intervention:** Everyone does not smoke.} \\
\texttt{2.  **Tar Deposit Distribution:** Since no one smokes, the probability of high tar deposit is the probability for nonsmokers, which is given as \$P(Tar=High | \textbackslash{}neg Smoking) = 0.77\$. Consequently, \$P(Tar=No | \textbackslash{}neg Smoking) = 1 {-} 0.77 = 0.23\$.} \\
\texttt{3.  **Lung Cancer Probability:**} \\
\texttt{    *   If Tar} \\
            }
        }
    }
    & & \\ \\

    \theutterance \stepcounter{utterance}  
    & & & \multicolumn{2}{p{0.3\linewidth}}{
        \cellcolor[rgb]{0.9,0.9,0.9}{
            \makecell[{{p{\linewidth}}}]{
                \texttt{\tiny{[GM$|$GM]}}
                \texttt{No} \\
            }
        }
    }
    & & \\ \\

    \theutterance \stepcounter{utterance}  
    & & & \multicolumn{2}{p{0.3\linewidth}}{
        \cellcolor[rgb]{0.9,0.9,0.9}{
            \makecell[{{p{\linewidth}}}]{
                \texttt{\tiny{[GM$|$GM]}}
                \texttt{no} \\
            }
        }
    }
    & & \\ \\

    \theutterance \stepcounter{utterance}  
    & & & \multicolumn{2}{p{0.3\linewidth}}{
        \cellcolor[rgb]{0.9,0.9,0.9}{
            \makecell[{{p{\linewidth}}}]{
                \texttt{\tiny{[GM$|$GM]}}
                \texttt{game\_result = WIN} \\
            }
        }
    }
    & & \\ \\

\end{supertabular}
}

\end{document}
